\documentclass[12pt]{article}
\usepackage{amsmath,amssymb,amsfonts}
\usepackage[margin=1in]{geometry}

\begin{document}

\section*{Chapter 2, Problem 1}

\subsection*{Problem Statement}
\textit{The shortest line (polar coordinates).} Find the shortest distance between two points using polar coordinates, i.e., using them as a line element:
\[
  d\ell^2 = dr^2 + r^2\,d\theta^2.
\]

\subsection*{Solution}

In polar coordinates, the arc length element is
\[
  d\ell = \sqrt{dr^2 + r^2\,d\theta^2}.
\]
We choose $\theta$ as the independent variable, so that $r = r(\theta)$ and
\[
  d\ell = \sqrt{r'^2 + r^2}\,d\theta,
\]
where $r' = dr/d\theta$. The total length between two points is
\[
  L = \int_{\theta_1}^{\theta_2} \sqrt{r'^2 + r^2}\,d\theta.
\]
The integrand is $f(r, r') = \sqrt{r'^2 + r^2}$, which does not depend explicitly on $\theta$. By the first integral (Eq.~2.4 of the text), we have
\[
  r'\,\frac{\partial f}{\partial r'} - f = \text{const}.
\]
We compute:
\[
  \frac{\partial f}{\partial r'} = \frac{r'}{\sqrt{r'^2 + r^2}}.
\]
Therefore,
\[
  \frac{r'^2}{\sqrt{r'^2 + r^2}} - \sqrt{r'^2 + r^2} = -C,
\]
where $C$ is a positive constant. Simplifying:
\[
  \frac{r'^2 - r'^2 - r^2}{\sqrt{r'^2 + r^2}} = -C
  \quad \Longrightarrow \quad
  \frac{r^2}{\sqrt{r'^2 + r^2}} = C.
\]
Squaring both sides:
\[
  r^4 = C^2(r'^2 + r^2)
  \quad \Longrightarrow \quad
  r'^2 = \frac{r^4 - C^2 r^2}{C^2} = \frac{r^2(r^2 - C^2)}{C^2}.
\]
Taking the square root and separating variables:
\[
  \frac{dr}{r\sqrt{r^2 - C^2}} = \frac{d\theta}{C}.
\]
We substitute $r = C\sec\phi$, so $dr = C\sec\phi\tan\phi\,d\phi$ and $\sqrt{r^2-C^2} = C\tan\phi$:
\[
  \frac{C\sec\phi\tan\phi\,d\phi}{C\sec\phi \cdot C\tan\phi} = \frac{d\theta}{C}
  \quad \Longrightarrow \quad
  \frac{d\phi}{C} = \frac{d\theta}{C}
  \quad \Longrightarrow \quad
  \phi = \theta - \alpha,
\]
where $\alpha$ is a constant of integration. Since $r = C\sec\phi$:
\[
  r = \frac{C}{\cos(\theta - \alpha)}
  \quad \Longrightarrow \quad
  r\cos(\theta - \alpha) = C.
\]
Expanding:
\[
  r\cos\theta\cos\alpha + r\sin\theta\sin\alpha = C.
\]
Using $x = r\cos\theta$ and $y = r\sin\theta$:
\[
  x\cos\alpha + y\sin\alpha = C,
\]
which is the equation of a straight line in Cartesian coordinates.

Thus, the shortest distance between two points in polar coordinates is indeed a straight line, as expected.

\end{document}
